\documentclass{BeamerTemplate}

\title{Module 4 - Loi normale et théorème central limite}

\subtitle{STT-1920 Méthodes statistiques}
\date{\today}

\begin{document}
\begin{frame}[t]
  \titlepage
\end{frame}

%%%%%%%%%%%%%%%%%%%%%%%%%%%%%%%%%%%
\section{Variables aléatoires continues}
%%%%%%%%%%%%%%%%%%%%%%%%%%%%%%%%%%%

\begin{frame}[t]{Variables aléatoires continues}

\begin{Boite}{Variable aléatoire continue}
Variable qui peut prendre toutes les valeurs dans un intervalle (ou $\mathbb{R}$).\\
La probabilité d'une valeur exacte est \alert{nulle} : $P(X = x) = 0$.
\end{Boite}

\medskip
\begin{Boite}{Densité de probabilité $f(x)$}
Fonction telle que :
$$P(a \le X \le b) = \int_a^b f(x)\,dx, \qquad f(x) \ge 0, \quad \int_{-\infty}^{+\infty} f(x)\,dx = 1$$
\end{Boite}

\medskip
L'espérance et la variance sont définies comme pour les lois discrètes (sommes $\to$ intégrales).

\end{frame}

%%%%%%%%%%%%%%%%%%%%%%%%%%%%%%%%%%%
\section{Loi normale}
%%%%%%%%%%%%%%%%%%%%%%%%%%%%%%%%%%%

\begin{frame}[t]{Loi normale}

\begin{Boite}{$X \sim \mathcal{N}(\mu, \sigma^2)$}
$$f(x) = \frac{1}{\sigma\sqrt{2\pi}}\exp\!\left(-\frac{(x-\mu)^2}{2\sigma^2}\right), \quad x \in \mathbb{R}$$
$$E(X) = \mu \qquad V(X) = \sigma^2$$
\end{Boite}

\medskip
Propriétés :
\begin{itemize}\itemsep3mm
  \item Courbe en forme de \alert{cloche}, symétrique autour de $\mu$.
  \item Paramètre $\mu$ : centre (position).
  \item Paramètre $\sigma$ : largeur (dispersion).
\end{itemize}

\end{frame}

\begin{frame}[t]{Règle empirique (68-95-99,7)}

Pour $X \sim \mathcal{N}(\mu, \sigma^2)$ :

\bigskip
\begin{itemize}\itemsep6mm
  \item $P(\mu - \sigma \le X \le \mu + \sigma) \approx 0{,}68$
  \item $P(\mu - 2\sigma \le X \le \mu + 2\sigma) \approx 0{,}95$
  \item $P(\mu - 3\sigma \le X \le \mu + 3\sigma) \approx 0{,}997$
\end{itemize}

\bigskip
$\Rightarrow$ Presque toutes les observations se trouvent à moins de 3 écarts-types de la moyenne.

\end{frame}

%%%%%%%%%%%%%%%%%%%%%%%%%%%%%%%%%%%
\section{Standardisation}
%%%%%%%%%%%%%%%%%%%%%%%%%%%%%%%%%%%

\begin{frame}[t]{Loi normale centrée réduite}

\begin{Boite}{Standardisation}
Si $X \sim \mathcal{N}(\mu, \sigma^2)$, alors :
$$Z = \frac{X - \mu}{\sigma} \sim \mathcal{N}(0, 1)$$
\end{Boite}

\begin{Boite}{Calcul de probabilité}
$$P(X \le x) = P\!\left(Z \le \frac{x-\mu}{\sigma}\right) = \Phi\!\left(\frac{x-\mu}{\sigma}\right)$$
\end{Boite}

\medskip
$\Phi(z)$ est la fonction de répartition de $\mathcal{N}(0,1)$ (table de la loi normale).

\end{frame}

\begin{frame}[t]{Calcul de probabilités — exemples}

$X \sim \mathcal{N}(8, 25)$ (moyenne $\mu = 8$, variance $\sigma^2 = 25$, donc $\sigma = 5$).

\bigskip
\begin{itemize}\itemsep6mm
  \item $P(X \le 13) = \Phi\!\left(\dfrac{13-8}{5}\right) = \Phi(1) \approx \alert{0{,}841}$
  \item $P(X > 3) = 1 - \Phi\!\left(\dfrac{3-8}{5}\right) = 1 - \Phi(-1) = \Phi(1) \approx 0{,}841$
  \item $P(3 \le X \le 13) = \Phi(1) - \Phi(-1) \approx 0{,}683$
\end{itemize}

\end{frame}

\begin{frame}[t]{Quantiles de la loi normale}

On cherche parfois $x$ tel que $P(X \le x) = p$.

\bigskip
\begin{Boite}{Quantile de $\mathcal{N}(\mu, \sigma^2)$}
$$x = \mu + z_p \cdot \sigma$$
où $z_p = \Phi^{-1}(p)$ est le quantile d'ordre $p$ de $\mathcal{N}(0,1)$.
\end{Boite}

\medskip
Quantiles usuels de $\mathcal{N}(0,1)$ :
\begin{center}
\begin{tabular}{|c|c|c|c|c|c|}
\hline
$p$ & 0{,}90 & 0{,}95 & 0{,}975 & 0{,}99 & 0{,}995 \\ \hline
$z_p$ & 1{,}282 & 1{,}645 & 1{,}960 & 2{,}326 & 2{,}576 \\ \hline
\end{tabular}
\end{center}

\end{frame}

%%%%%%%%%%%%%%%%%%%%%%%%%%%%%%%%%%%
\section{Théorème central limite}
%%%%%%%%%%%%%%%%%%%%%%%%%%%%%%%%%%%

\begin{frame}[t]{Théorème central limite (TCL)}

\begin{Boite}{Théorème central limite}
Soient $X_1, X_2, \ldots, X_n$ des v.a. i.i.d. d'espérance $\mu$ et de variance $\sigma^2 < \infty$.\\
Quand $n \to \infty$ :
$$\overline{X} = \frac{1}{n}\sum_{i=1}^{n} X_i \xrightarrow{\text{loi}} \mathcal{N}\!\left(\mu, \frac{\sigma^2}{n}\right)$$
\end{Boite}

\medskip
\begin{itemize}\itemsep4mm
  \item \alert{Peu importe} la distribution des $X_i$, la moyenne échantillonnale est approximativement normale pour $n$ suffisamment grand.
  \item En pratique : $n \ge 30$ suffit souvent.
\end{itemize}

\end{frame}

\begin{frame}[t]{Pourquoi le TCL est fondamental}

Le TCL justifie l'utilisation de la loi normale pour l'inférence statistique, même quand la population n'est pas normale.

\bigskip
\begin{columns}[t]
\column{0.5\textwidth}
\textbf{Ce qu'on sait :}
\begin{itemize}\itemsep3mm
  \item $X_i$ i.i.d., espérance $\mu$, variance $\sigma^2$
\end{itemize}

\column{0.5\textwidth}
\textbf{Ce qu'on peut dire :}
\begin{itemize}\itemsep3mm
  \item $E(\overline{X}) = \mu$
  \item $V(\overline{X}) = \sigma^2/n$
  \item $\overline{X} \approx \mathcal{N}(\mu, \sigma^2/n)$ pour $n$ grand
\end{itemize}
\end{columns}

\bigskip
$\Rightarrow$ Toute la théorie des intervalles de confiance et des tests d'hypothèses repose sur ce résultat.

\end{frame}

\begin{frame}[t]{TCL — illustration}

\underline{Exemple :} Une infirmière note le nombre quotidien de naissances ($N \sim \text{Poisson}(8)$) sur 365 jours.

\bigskip
\begin{itemize}\itemsep5mm
  \item $E(N) = 8$, $V(N) = 8$.
  \item Après 365 jours, par le TCL :
  $$\overline{N} \approx \mathcal{N}\!\left(8, \frac{8}{365}\right) = \mathcal{N}(8, 0{,}022)$$
  \item La moyenne observée sera très proche de 8.
\end{itemize}

\end{frame}

\end{document}
